\chapter{Methodik und Entwicklungsablauf}

\section{Grundprinzip}

Die Entwicklung folgt einem inkrementellen MBSE-Prozess. Zuerst werden Zweck, Kontext und Systemgrenze verstanden. Danach werden Stakeholder-Beduerfnisse in testbare Anforderungen ueberfuehrt. Aus Anforderungen entstehen Use Cases, Funktionen, logische und physische Architekturen, Schnittstellen, \ros-Elemente und Verifikationsfaelle.

Das Lehrbuch beschreibt MBSE nicht als Produktion moeglichst vieler Diagramme, sondern als systematische Verknuepfung von Informationen: Stakeholder-Bedarf, Anforderung, Funktion, Systemelement und Testfall muessen nachvollziehbar miteinander verbunden sein \cite[Kap.~2]{becker2026sysml}. Diese Traceability ist fuer das Roboterprojekt besonders wichtig, weil Aenderungen an Navigation, Safety, Sensorik oder Energieversorgung mehrere Entwicklungsbereiche gleichzeitig betreffen koennen.

\section{Iterativer Ablauf}

Der Entwicklungsablauf wird in Iterationen organisiert. Jede Iteration erzeugt konkrete Artefakte und aktualisiert dieses Dokument.

\begin{figure}[htbp]
  \centering
  \resizebox{\textwidth}{!}{\begin{tikzpicture}[
  node distance=9mm,
  step/.style={draw, rounded corners=2pt, align=center, minimum width=31mm, minimum height=8mm, fill=gray!8},
  artifact/.style={draw, rounded corners=2pt, align=center, minimum width=31mm, minimum height=8mm, fill=blue!6},
  arrow/.style={-{Latex[length=2mm]}, thick},
  feedback/.style={-{Latex[length=2mm]}, thick, dashed}
]
\node[step] (goal) {Systemzweck\\und Kontext};
\node[step, right=of goal] (needs) {Stakeholder\\Needs};
\node[step, right=of needs] (req) {Requirements\\und Use Cases};
\node[step, right=of req] (func) {Funktionen\\und Verhalten};
\node[step, below=of func] (arch) {Architektur\\und Schnittstellen};
\node[step, left=of arch] (mvp) {MVP-Schnitt\\im Zielsystem};
\node[step, left=of mvp] (impl) {HW/SW, ROS2\\und Mechanik};
\node[artifact, left=of impl] (verify) {Verifikation\\und Traceability};

\draw[arrow] (goal) -- (needs);
\draw[arrow] (needs) -- (req);
\draw[arrow] (req) -- (func);
\draw[arrow] (func) -- (arch);
\draw[arrow] (arch) -- (mvp);
\draw[arrow] (mvp) -- (impl);
\draw[arrow] (impl) -- (verify);
\draw[feedback] (verify.north) to[out=90,in=270] node[above, align=center, font=\small] {Erkenntnisse aktualisieren\\Modell und Entscheidungen} (needs.south);
\end{tikzpicture}
}
  \caption{Inkrementeller Entwicklungsablauf vom Systemzweck bis zur Verifikation}
  \label{fig:development-flow}
\end{figure}

\begin{longtable}{L{0.16\textwidth} L{0.30\textwidth} L{0.42\textwidth}}
\toprule
Iteration & Schwerpunkt & Ergebnisartefakte \\
\midrule
A & Systemdefinition & Zielsystem, Systemgrenze, Einsatzumgebung, Stakeholder, Annahmen, erste Risiken \\
B & Requirements und Use Cases & Stakeholder Needs, Requirements, Use Cases, Akzeptanzkriterien \\
C & MVP-Definition & MVP-Scope, Nicht-MVP, Erfolgskriterien, MVP-Traceability \\
D & Systemarchitektur & Funktionen, logische Architektur, physische Architektur, Schnittstellen, Betriebsmodi \\
E & Implementierungsarchitektur & HW/SW-Allokation, Rechnerarchitektur, Sensorik, Elektronik, \ros-Mapping \\
F & Verifikation & Testfaelle, Simulation, Integrationsstrategie, Systemtests, Traceability-Matrix \\
\bottomrule
\end{longtable}

\section{Single Source of Truth}

Das Gesamtmodell bleibt die fachliche Referenz. Der MVP wird nicht als separates Modell erstellt, sondern als markierter Ausschnitt des Zielsystems. Jedes relevante Element erhaelt daher mindestens:
\begin{itemize}
  \item eine stabile ID,
  \item einen Status: \emph{MVP}, \emph{Release 1}, \emph{Future} oder \emph{Offen},
  \item eine Prioritaet: \emph{Must}, \emph{Should}, \emph{Could},
  \item eine Quelle oder Begruendung,
  \item eine geplante Nachweismethode.
\end{itemize}

\section{Rolle von KI im Prozess}

KI wird als Engineering-Partner eingesetzt, nicht als alleinige Freigabeinstanz. Sie unterstuetzt beim Strukturieren, Generieren und Pruefen von Artefakten. Formale Modellpruefung, Berechnungen, Simulation, Compiler, Tests, CAD-/EDA-Pruefungen und menschliche Reviews bleiben entscheidend.

Die Projektregel lautet:
\begin{quote}
KI darf Entwicklungsartefakte vorschlagen und konsistent halten. Technische Gueltigkeit entsteht erst durch nachvollziehbare Anforderungen, begruendete Entscheidungen, formale Modellierung, Berechnung, Simulation, Test und Review.
\end{quote}

\section{Qualitaetskriterien fuer Artefakte}

Ein Artefakt gilt erst dann als belastbar, wenn folgende Fragen beantwortet sind:
\begin{itemize}
  \item Welche Anforderung oder welches Stakeholder-Beduerfnis begruendet das Artefakt?
  \item Welche Alternative wurde betrachtet?
  \item Warum wurde die aktuelle Loesung gewaehlt?
  \item Welche Schnittstellen, Tests und Risiken entstehen daraus?
  \item Wie wirkt sich eine spaetere Aenderung auf andere Artefakte aus?
\end{itemize}

\section{Entwicklungsprinzipien aus Iteration A}

Aus der Vervollstaendigung von Iteration A werden folgende Entwicklungsprinzipien als Projektregeln uebernommen:

\begin{longtable}{L{0.14\textwidth} L{0.28\textwidth} L{0.44\textwidth}}
\toprule
ID & Prinzip & Bedeutung \\
\midrule
DP-01 & Requirement vor Loesung & Technologien werden erst festgelegt, wenn Anforderungen sie begruenden. \\
DP-02 & Gesamtmodell vor MVP & Das Gesamtprojekt wird modelliert; der MVP ist eine Teilmenge davon. \\
DP-03 & Keine Wegwerfarchitektur & MVP-Komponenten sollen moeglichst in die spaetere Architektur ueberfuehrbar sein. \\
DP-04 & Traceability & Needs, Requirements, Use Cases, Funktionen, Architektur, Implementierung und Verifikation werden verknuepft. \\
DP-05 & Entscheidungen dokumentieren & Groessere Architekturentscheidungen erhalten einen Decision Record. \\
DP-06 & Unsicherheit sichtbar machen & Requirements, Decisions, Assumptions, Open Issues und Risks werden getrennt dokumentiert. \\
DP-07 & Schrittweise Verifikation & Nicht erst der fertige Roboter wird getestet; Teilmodelle, Simulationen, Baugruppen und Software werden frueh geprueft. \\
\bottomrule
\end{longtable}

\section{Artefakttypen und Umgang mit Unsicherheit}

Fuer die weitere Entwicklung wird strikt unterschieden zwischen:
\begin{itemize}
  \item \emph{Requirement}: nachzuweisende Forderung an das System,
  \item \emph{Decision}: getroffene und begruendete Architektur- oder Prozessentscheidung,
  \item \emph{Assumption}: vorlaeufige Annahme, die spaeter bestaetigt oder ersetzt werden muss,
  \item \emph{Open Issue}: offene Frage mit Auswirkung auf Anforderungen, Architektur oder Tests,
  \item \emph{Risk}: moegliches Ereignis oder technische Unsicherheit mit negativer Auswirkung.
\end{itemize}

Diese Trennung ist wichtig, damit offene Fragen nicht versehentlich als Anforderungen erscheinen und Loesungsideen nicht zu frueh als feste Randbedingungen behandelt werden.
